\section{Corrections to the v0.4 study}
\label{sec:corrections}

The v0.4 study \cite{fishbot04} and its supporting artifacts contain claims that
the measurements of \S\ref{sec:diagnosis} contradict. They are stated here as
corrections, in one place, rather than quietly patched. Each entry gives what was
claimed, where, what is now measured, and what replaces it. Where the v0.4
\emph{paper} and a v0.4 \emph{artifact} disagree with each other, that is said
explicitly, because in the load-bearing case the paper is right and the artifact
is wrong.

\paragraph{C1. The information-freeze account of non-termination is withdrawn.}
The v0.4 termination artifact, \filepath{research/v04/results/E11-termination.md},
states that the locked-half-suit theorem implies that such a half-suit is frozen
``and, for the same reason, that no further information about its allocation can
ever arrive,'' and concludes that two policies which both decline to cash an
uncertain half-suit therefore deadlock. It repeats the claim as ``Theorem~1
freezes information as well as ownership,'' and closes by offering that
proposition as a contribution to the rules literature. That is false, and all
three statements are retracted.

It is false at the level of the rules before any measurement: an ask inside a
half-suit locked to the asker's own team satisfies every legality condition, and
the deduction machinery publishes the asker's own exclusion whether the ask
succeeds or not. It is false as measured: \numCertRate\% of such asks strictly
raise a teammate's exact probability of the correct allocation, mean
$+\numCertMean$, maximum $+\numCertMax$, and \numEntropyRate\% strictly lower the
exact marginal entropy, by a mean of \numEntropyMean{} nats
(\S\ref{sec:diag-information}). And it is not the mechanism of the observed
deadlock in any case: in \numNoLockPct\% of long games no half-suit is locked to
any team at any point in the dead run, and only \numCycleInLock\% of long games
have their entire cycle inside a lock. What replaces it is
\S\ref{sec:diag-cycle}: the deadlock is a deterministic two-question cycle at a
fixed point of the public information state, produced by an ask score in which
ungated ownership features outbid the hit-probability term at $p = 0$.

The v0.4 \emph{paper} is not implicated in this error and should not be read as
if it were. Its Observation on locked half-suits states the correct position ---
that ownership freezes while allocation information need not, that asks inside a
locked half-suit stay legal and carry certificates, and that ownership
monotonicity alone establishes neither informational freezing nor
non-termination. The error is confined to the artifact, was flagged in the
engine's own specification document, and was never corrected there.

\paragraph{C2. The reported termination property was purchased, and the price was
not reported.} The v0.4 study reports an action-limit rate of \numVfourLimitRate{}
across its experiments and describes the graduated event-count escalation as the
mechanism. Both statements are accurate. What was not reported is what the
escalation does in the population where it actually fires. In mirror play
\numPostHorizonDecl{} declarations are made at or after the forcing horizon and
\numPostHorizonWrong\% of them are wrong, against a baseline declaration error of
\numDeclWrong\%; bucketed by the confidence the policy itself attaches,
\numLowConfDecl{} declarations in \numLowConfGames{} games are made below a
self-assessed probability of \numLowConfThresh{} and every one is wrong. The escalation's second stage returns ``declare'' before
inspecting the allocation probability at all, and removing it is worth
\numStageTwoCost{} points against a mirror opponent while changing nothing
against any weaker style tested. A zero cap-hit rate against a panel of weaker
policies is therefore not evidence that termination was solved; it is evidence
that the horizon was rarely reached in that panel. The correct statement is that
\vfour{} converts non-termination into misdeclaration, and that the conversion
rate is a strong-opponent tax.

\paragraph{C3. Forced-endgame declarations were reported as inaccurate and
recorded as a known gap; they are a certainty of loss with a mechanical cause.}
The v0.4 evaluation tables carry a forced-declaration accuracy of zero, and the
engine specification records that the allocation routine takes a per-card argmax
with no capacity constraint. The two were never connected. \S\ref{sec:diag-forced}
connects them: over \numForcedGames{} distinct mirror games, all \numForcedN{}
forced declarations violate the declarer's own capacity constraint, all have
exact posterior probability zero, and the declarer's own stated confidence is
exactly zero in every one. The best feasible allocation would have been correct
in \numFeasibleRight\% of them. This is not a hard inference problem being lost;
it is a policy naming an allocation it has itself scored at zero. Two notes
attach to the same measurement, neither of which touches the rate. The count of
\numForcedN{} corrects, downward by half, the figure the diagnosis run first
reported: in a mirror match the six seat rotations are three deal shifts times
two team orientations, both orientations play byte-identical games, and every
such run splits its forced declarations exactly evenly between them. And the
result does not rest on that run alone --- an instrument sharing no code with it,
recomputing each declaration from the opening deal, reproduces it at five further
mirror seeds.

\paragraph{C4. The known gap about the value-function coefficients has the sign
backwards.} The v0.4 specification records that the sixteen compiled value
coefficients are not the ones in the fitting artifact, framed as a defect. The
compiled vector is the better of the two: substituting the fitted vector costs
\numEfourteenRange{} win-rate points, including \numEfourteenCost{} points
against \vthree{}, with the sign stable across nine opponents and three
independent seed banks. Two riders on that figure are corrections to the present
study rather than to the v0.4 one. The size is not uniform across the panel: at a
fresh seed the mirror is the \emph{smallest} effect and the only one whose
cluster-bootstrap interval includes zero, so a narrower per-opponent range stated
in an earlier draft of this analysis, and the reading that the compiled vector's
advantage is largest in mirror play, are both withdrawn. The real defect is
larger and different. The model is
effectively one feature: held-out $R^2$ is \numValueRSq{} for all sixteen and
\numValueRSqScore{} for a bias term plus the score differential alone, two
features are algebraically identical, and five more cancel at every call site ---
moving all five to $\pm\numDeadFeatureSwing$ changes none of \numDeadFeatureDecisions{} declaration
decisions. The v0.4 limitation that no goodness-of-fit figure was reported for
the deployed coefficients is upgraded here to a measurement, and it is not
flattering.

\paragraph{C5. The policy prior is one channel, not two.} The v0.4 study
describes the prior as two coefficients, a weight on having asked in a half-suit
and a weight on having taken turns without asking there, and reports a
frozen-policy ablation for the pair. The second is not an independent channel:
the exponent rearranges so that its term is card-independent, and the capacity
normalisation then removes it. Deleting it costs nothing measurable ---
\numPhiWin\% [\numPhiCI] against the honest mirror, a paired delta of
\numPhiPaired{} points across a four-opponent panel, and a fifteen-fold sweep flat
inside its intervals. The ablation result for the pair stands, and the
attribution of it to two channels does not. This also corrects a hypothesis of
the present study rather than of the v0.4 one: the silence channel that the
project owner's manoeuvre attacks was not misreading silence, it was not reading
it at all.

\paragraph{C6. Asks inside a locked half-suit: both halves of the v0.4 hedge are
now measured, and they point in opposite directions.} The v0.4 limitations
describe these asks as ``a channel the fitted configuration has no term for
\dots\ unexploited here, not shown to be worthless.'' They are not worthless:
\S\ref{sec:diag-information} prices them, and a greedy ladder of them reaches
certainty in \numLadderRate\% of attempts at an information exchange rate of
about $\numLeakRatio\!:\!1$ in the sender's favour. They are also not usable in
the form the hedge implies. Selecting them requires the actor to know its team
owns the half-suit, and that condition is provable from public information at
only \numProvableOwn\% of mirror decisions; at the deadlock states the exact
probability an owning team assigns to its own ownership averages
\numOwnProbMean{} and peaks at \numOwnProbMax{}, so not one of the
\numOwnProbPairs{} pairs measured comes near the \numProvableThresh{} bar. The channel has to be priced against ownership probability rather
than gated on ownership. The first pass of the present study wrote ``provably
owns'' where it meant ``in fact owns,'' and that wording is corrected here.

\paragraph{C7. The eight-level forced-endgame willingness ladder does not
function, and leaks more than the rules license.} The v0.4 dialect table presents
the ladder as a graded willingness mechanism and contrasts it with v0.3's
two-level version. It is inert: the statistic it thresholds is exactly zero in
\numRungZeroPairs{} of \numRungZeroPairs{} surveyed (player, half-suit) pairs, a
consequence of C3, and four ladder shapes including one with no rungs at all
produce bit-identical output. Separately, a ladder of this shape transmits
\numLadderLeak{} bits per candidate, against the roughly one bit of willingness
the rules make public. M2 restores the ladder's function for free; the leak
figure is the reason \S\ref{sec:mech-rejected} declines to extend the same device
to mid-game turn transfer, where the audience still holds cards.

\paragraph{C8. Waiting was free by construction, so observing that waiting
dominates proves nothing.} The v0.4 declaration rule compares the value of acting
against the value of waiting, and the waiting branch is evaluated by passing
all-zero deltas to a value model in which the public event count appears nowhere.
The value of waiting at time $t$ therefore equals the value of waiting at $t+1$
as an identity of the implementation, not as an approximation or a finding. Any
argument that patience weakly dominates in that model is self-fulfilling, and the
only thing that ever broke the tie was the hand-set horizon whose cost C2
records. \S\ref{sec:mech-rejected} reports that the literature's remedy for this
--- a strictly positive holding cost --- was implemented and swept, and that it
shifts the mean without touching the tail, so the correction here is not that
v0.4 chose the wrong cost but that the well-posedness of its stopping problem
came from a move cap rather than from anything in the model.

\paragraph{C9. Scope: the published head-to-head result stands; its
generalisation caveat now has an instance.} \vfour{}'s reported
\numVfourVsVthree\% against \vthree{} is not contradicted by anything here, and
the v0.4 study was explicit that it established nothing about generalisation to
unseen opponent populations. \S\ref{sec:diag-baseline} supplies a concrete
instance of that caveat biting: on the same instrument, against a population the
panel did not contain, the dead-ask rate differs by a factor of fourteen and a
third of games contain an extended run of guaranteed misses. The lesson is about
evaluation design and is acted on in M10 --- the mirror belongs in the fitting
panel and the pathology statistics belong in the commit gate.
